\chapter{Power and Gas Trading}\label{ch:s2:power-and-gas-trading}

Power bought the day before for delivery tomorrow and power bought an hour before delivery are different products. The day-ahead auction clears on
the wind and solar forecasts of the previous morning; the intraday market prices what actually blows and shines. Kiesel and Paraschiv found that German
intraday prices adjust to forecast errors in renewables. On Germany's 2024--2025 day-ahead prices, with intraday prices simulated from planted wind
forecast errors, a 1~MW, 2~MWh battery scheduled on the day-ahead auction earns \euro69\,900 a year. Re-optimising it against intraday prices adds
\euro2\,900. A trader who forecasts half the wind error earns \euro2.54 per MWh traded between the two markets. The build is \texttt{firm.powergas}.

\section{Day-ahead against intraday}

\begin{definition}[Day-ahead-intraday spread]\label{def:s2:power-and-gas-trading:daid}
The \emph{day-ahead-intraday spread}\index{day-ahead-intraday spread} for a delivery hour is the day-ahead auction price less the intraday price for
the same hour; it is driven by what changed between the auction and delivery, above all the forecasts of wind, solar and demand.
\end{definition}

Germany's day-ahead prices tie tightly to residual load, the load left after wind and solar. Over 2024 and 2025 a regression of the hourly price on
residual load has a slope of \euro3.05 per MWh per GW and a correlation of 0.83. If the auction cleared on a wind forecast that was 1~GW too low, the
hour was priced for 1~GW more residual load than arrived, and the intraday price should be about \euro3 lower. \texttt{firm.powergas} builds intraday
prices this way: the day-ahead price less \euro3.05 per GW of unexpected wind, with forecast errors of 2~GW standard deviation that persist from hour to
hour (autocorrelation 0.9), plus \euro6 of other noise (\cref{lst:s2:power-and-gas-trading:trade}). These errors are planted: real intraday prices and
forecasts are licensed.

The evidence on real intraday prices is mixed in its details. Kiesel and Paraschiv found that intraday prices adjust asymmetrically to renewable
forecast errors and to trading volume, depending on a threshold in the demand expected to be covered by planned conventional capacity. Ziel, with a
regression model of German intraday prices, found no statistically significant evidence that positive and negative wind or solar errors have different
effects.

A trader whose forecast of the error has a correlation of 0.5 with it sells day-ahead and buys back intraday when expecting more wind than the auction
assumed, and does the reverse when expecting less; with a half-spread of \euro1 on the intraday side:

\begin{center}
\small
\begin{tabular}{@{}lccc@{}}
\toprule
forecast skill (correlation), half-spread & \euro{} per MWh traded & \euro{} a year per MW & IC\\
\midrule
0.5, \euro1 & 2.54 & 13\,600 & 0.36\\
0.3, \euro1 & 1.08 & 5\,800 & 0.21\\
0.5, \euro3 & 0.54 & 2\,900 & 0.36\\
\bottomrule
\end{tabular}
\end{center}

The trade is in the market in 61\% of hours. The edge is a forecast of the error, and it shrinks quickly with skill and with costs.

\section{Weather-driven positions}

\begin{definition}[Weather-driven position]\label{def:s2:power-and-gas-trading:weather}
A \emph{weather-driven position}\index{weather-driven position} is a position in power or gas taken on a forecast of weather (temperature, wind,
cloud, precipitation) that differs from the one in the market's prices, closed when the forecasts converge or the weather arrives.
\end{definition}

The day-ahead trade is the shortest \omterm{def:s2:power-and-gas-trading:weather}{weather-driven position}. Over days and weeks, temperature forecasts drive heating and cooling demand, and with
it gas and power prices. A trader who reads a forecast revision before the market prices it holds the same kind of edge, at a longer horizon and a
larger scale, with the same enemies: a forecast that the market already has, and costs. Book~3 (chapter~11) treats weather derivatives, which
hedge the weather itself.

\section{Batteries and flexible assets}

\begin{definition}[Battery arbitrage]\label{def:s2:power-and-gas-trading:battery}
\emph{Battery arbitrage}\index{battery arbitrage} is scheduling a battery to charge in cheap hours and discharge in dear ones, first against the
day-ahead auction and then by re-optimising against intraday prices, within its power, energy, efficiency and cycle limits.
\end{definition}

Germany's prices have the shape a battery needs (\cref{fig:s2:power-and-gas-trading:profile}): cheap at midday when solar peaks, dear in the evening.
Prices were negative in 457 hours in 2024 and 573 in 2025, down to $-$\euro250.32; at 13:00 local time 23.4\% of hours were negative, and at 19:00
none. The battery is scheduled each day by dynamic programming over its state of charge and the energy charged so far
(\cref{lst:s2:power-and-gas-trading:battery}), with a round-trip efficiency of 88\% and one full cycle a day. On the 727 days with 24 hours, knowing
each day's auction prices in advance, as in Book~3, it earns \euro69\,900 per MW a year. A quarter of its charging hours (24.8\%) have negative prices.

\begin{omfigure}
\begin{tikzpicture}
\begin{axis}[width=11cm,height=5.2cm,axis y line*=right,axis x line=none,ylabel={\small negative hours (\%)},
  tick label style={font=\footnotesize},xmin=-0.5,xmax=23.5,ymin=0,ymax=60,ybar,bar width=5pt,table/col sep=comma]
\addplot[fill=gray!30,draw=gray] table[x=hour,y=negative]{figdata/strategies-2/23-power-and-gas-trading/profile.csv};\label{plot:s2:powneg}
\end{axis}
\begin{axis}[width=11cm,height=5.2cm,axis y line*=left,xlabel={\small hour, German local time},ylabel={\small \euro/MWh, mean},
  tick label style={font=\footnotesize},xmin=-0.5,xmax=23.5,ymin=0,ymax=150,xtick={0,3,6,9,12,15,18,21},grid=major,grid style={gray!25},
  table/col sep=comma,legend style={font=\scriptsize,at={(0.5,-0.3)},anchor=north},legend columns=3]
\addplot[omThm,thick,mark=*,mark size=1pt] table[x=hour,y=p2024]{figdata/strategies-2/23-power-and-gas-trading/profile.csv};
\addlegendentry{2024}
\addplot[omDef,thick,mark=square*,mark size=1pt] table[x=hour,y=p2025]{figdata/strategies-2/23-power-and-gas-trading/profile.csv};
\addlegendentry{2025}
\addlegendimage{/pgfplots/refstyle=plot:s2:powneg}\addlegendentry{negative hours (right)}
\end{axis}
\end{tikzpicture}

\omcaption{German day-ahead prices (DE-LU) by hour of the day, averaged over 2024 and over 2025, and the share of hours with negative prices over both
years. Data: Bundesnetzagentur / SMARD.de (CC BY 4.0).}\label{fig:s2:power-and-gas-trading:profile}
\end{omfigure}

The battery that has sold its day-ahead schedule can change it intraday: buy back power it no longer wants to deliver, or sell more, paying the
half-spread on every MWh it changes. Re-optimising against the day's intraday prices, the schedule changes on 71\% of days:

\begin{center}
\small
\begin{tabular}{@{}lcc@{}}
\toprule
1~MW, 2~MWh, 88\% round trip & day-ahead (\euro{} a year) & intraday extra (\euro{} a year)\\
\midrule
one cycle a day, half-spread \euro1 & 69\,900 & 2\,900\\
one cycle, no spread & 69\,900 & 3\,800\\
one cycle, half-spread \euro3 & 69\,900 & 1\,600\\
two cycles a day, half-spread \euro1 & 83\,800 & 4\,900\\
\bottomrule
\end{tabular}
\end{center}

Both columns are generous. The day-ahead column knows the auction's prices before bidding into it; the intraday column re-optimises once with the
whole day's intraday prices known. The intraday extra is small because the planted errors move prices by a few euros, while the day's price range
averaged over a hundred (\cref{fig:s2:power-and-gas-trading:battery}).

\begin{omfigure}
\begin{tikzpicture}
\begin{axis}[width=11cm,height=5cm,xlabel={\small day (1 January 2024 to 31 December 2025, clock-change days omitted)},
  ylabel={\small cumulative \euro{} thousands per MW},tick label style={font=\footnotesize},xmin=0,xmax=730,ymin=0,ymax=155,
  grid=major,grid style={gray!25},table/col sep=comma,legend style={font=\scriptsize,at={(0.03,0.97)},anchor=north west}]
\addplot[omThm,thick] table[x=day,y=da]{figdata/strategies-2/23-power-and-gas-trading/battery.csv};
\addplot[omDef,thick,dashed] table[x=day,y=total]{figdata/strategies-2/23-power-and-gas-trading/battery.csv};
\legend{day-ahead schedule,with intraday re-optimisation}
\end{axis}
\end{tikzpicture}

\omcaption{Cumulative revenue of a 1~MW, 2~MWh battery on German day-ahead prices, 2024--2025, and with its schedule re-optimised against synthetic
intraday prices. Data: Bundesnetzagentur / SMARD.de (CC BY 4.0); \texttt{s2\_powergas.battery\_table}.}\label{fig:s2:power-and-gas-trading:battery}
\end{omfigure}

\section{Risk in power books}

Power cannot be stored in bulk, so every hour is its own product, and a power book is a book of 8\,760 contracts a year, each with its own weather.
Its risks are the forecast errors themselves, which are correlated across hours (a wrong wind forecast is wrong all afternoon); prices that jump
by hundreds of euros when a few gigawatts go missing (the 2024 maximum was \euro936.28); imbalance charges for delivering less than sold (Book~3,
chapter~6); and the limits of the physical asset, a battery that is full when prices go negative. A battery also degrades with each cycle, which the
table above does not charge. A second cycle a day adds \euro13\,800 of day-ahead revenue at a cost in battery life that the owner must price.

\section{Strategy files}

\begin{strategyfile}{Forecast-error intraday trade}\label{strat:s2:power-and-gas-trading:forecast}
\sfield{Who pays you, and why} Generators and suppliers who must balance their forecast errors intraday.
\sfield{Instruments and venues} Day-ahead auction; continuous intraday market.
\sfield{Signal} The trader's forecast of wind, solar and load against the auction's.
\sfield{Sizing and execution} Positions scaled to the forecast's confidence; closed intraday.
\sfield{Costs} Intraday bid-ask; imbalance charges on what is not closed.
\sfield{How it dies} Better public forecasts; costs.
\sfield{Horizon, capacity, infrastructure} Hours; weather models and fast intraday access.
\sfield{Backtest honestly} Forecasts as of the auction, not later ones.
\sfield{Sources} Kiesel and Paraschiv (2017); Ziel (2017); this chapter: \euro2.54 per MWh at a skill of 0.5.
\end{strategyfile}

\begin{strategyfile}{Weather-driven gas position}\label{strat:s2:power-and-gas-trading:weather}
\sfield{Who pays you, and why} Market participants who price yesterday's forecast.
\sfield{Instruments and venues} Gas futures for the next weeks; power baseload.
\sfield{Signal} Revisions of heating and cooling degree-day forecasts against the priced forecast.
\sfield{Sizing and execution} In proportion to the revision's expected demand effect.
\sfield{Costs} Futures bid-ask.
\sfield{How it dies} Everyone has the same models; revisions priced in minutes.
\sfield{Horizon, capacity, infrastructure} Days to weeks.
\sfield{Backtest honestly} Forecast vintages as issued.
\sfield{Sources} No performance figure verified.
\end{strategyfile}

\begin{strategyfile}{Battery day-ahead arbitrage}\label{strat:s2:power-and-gas-trading:battery}
\sfield{Who pays you, and why} The daily price shape: solar at midday, demand in the evening.
\sfield{Instruments and venues} A battery; day-ahead auction; intraday market; reserves.
\sfield{Signal} Forecast day-ahead prices; then intraday prices.
\sfield{Sizing and execution} A daily schedule by dynamic programming; re-optimised intraday.
\sfield{Costs} Round-trip losses; degradation; spreads.
\sfield{How it dies} More batteries flattening the shape.
\sfield{Horizon, capacity, infrastructure} Hours; the asset.
\sfield{Backtest honestly} Price forecasts, not realised auction prices; degradation per cycle.
\sfield{Sources} This chapter: \euro69\,900 per MW a year day-ahead with foresight, \euro2\,900 more intraday.
\end{strategyfile}

\begin{strategyfile}{Spark-spread plant optimisation}\label{strat:s2:power-and-gas-trading:plant}
\sfield{Who pays you, and why} The spread between power and fuel in the hours a plant runs.
\sfield{Instruments and venues} A gas plant or a tolling contract; power, gas and carbon forwards.
\sfield{Signal} The clean spark spread by hour against the plant's costs.
\sfield{Sizing and execution} Dispatch when the spread exceeds costs; hedge forward in blocks.
\sfield{Costs} Start-up costs; minimum run times.
\sfield{How it dies} Renewables that take the plant's hours.
\sfield{Horizon, capacity, infrastructure} Hours to years.
\sfield{Backtest honestly} Start-up and ramping constraints.
\sfield{Sources} No performance figure verified.
\end{strategyfile}

\begin{strategyfile}{Negative-price hours}\label{strat:s2:power-and-gas-trading:negative}
\sfield{Who pays you, and why} Producers who must or will run at any price (subsidy schemes, inflexible plants).
\sfield{Instruments and venues} Flexible load or storage; day-ahead and intraday markets.
\sfield{Signal} Forecast solar and wind against load at midday and on holidays.
\sfield{Sizing and execution} Consume or charge in negative hours.
\sfield{Costs} Network charges on consumption.
\sfield{How it dies} Subsidy rules that stop paying in negative hours; more flexible demand.
\sfield{Horizon, capacity, infrastructure} Hours.
\sfield{Backtest honestly} Network charges; the forecast, not the realised price.
\sfield{Sources} SMARD data: 457 negative hours in 2024 and 573 in 2025.
\end{strategyfile}

\section{Tutorial: the wind forecast}\label{tut:s2:power-and-gas-trading:tut}

\begin{tutorial}
\textbf{Goal.} Measure the German price profile and negative hours, simulate intraday prices from wind forecast errors, trade the errors, and schedule a
battery day-ahead and then intraday. \textbf{End state:} the three tables and two figures.
\begin{tutsteps}
\item \textbf{Intraday prices and the forecast trade}.
\omcode{code/firm/powergas/firm_powergas.py}{42}{63}{Wind errors move intraday prices; a partial forecast trades them.}\label{lst:s2:power-and-gas-trading:trade}
\item \textbf{The battery}.
\omcode{code/firm/powergas/firm_powergas.py}{66}{101}{Daily dynamic programming, day-ahead or against a committed schedule.}\label{lst:s2:power-and-gas-trading:battery}
\item \textbf{Run} \texttt{real()}, \texttt{trade()}, \texttt{battery\_table()} and \texttt{fig\_powergas.py}.
\end{tutsteps}
\textbf{What to change next.} Bid the battery on a price forecast instead of the auction's prices; re-optimise intraday hour by hour; add a degradation
cost per cycle.
\end{tutorial}

\section{Build: power and gas}\label{bld:s2:power-and-gas-trading:powergas}

\begin{build}
\bfield{Purpose} Intraday prices from forecast errors, the forecast-error trade, and battery scheduling day-ahead and intraday.
\bfield{Interface} \texttt{PowerConfig(\dots)}, \texttt{intraday(da, cfg)}, \texttt{forecast\_trade(da, idp, err, cfg)}, \texttt{battery(prices,
cfg, committed)}.
\bfield{Rules} Charging costs $1/\sqrt{\eta}$ of the grid energy, discharging returns $\sqrt{\eta}$; the battery ends each day empty; deviations
from a committed schedule pay the half-spread.
\bfield{Acceptance tests} \texttt{code/firm/powergas/tests/}: no errors, no spread; a perfect forecast never loses without costs; the battery by hand;
re-optimising on the committed prices changes nothing.
\bfield{Stretch} Quarter-hour products; reserves; degradation.
\end{build}

\begin{omsources}
\item R.~Kiesel and F.~Paraschiv, ``Econometric analysis of 15-minute intraday electricity prices'', \emph{Energy Economics} 64, 2017.
\item F.~Ziel, ``Modeling the impact of wind and solar power forecasting errors on intraday electricity prices'', \emph{14th International Conference
on the European Energy Market}, 2017.
\item Bundesnetzagentur / SMARD.de, hourly day-ahead prices, generation and load for Germany, 2024--2025, CC BY 4.0 (Book~3's file).
\end{omsources}

\section{Exercises}

\begin{exercise}[$\star$]\label{exo:s2:power-and-gas-trading:1}
Wind comes in 1.5~GW above the auction's forecast. At \euro3.05 per MWh per GW, how far should the intraday price be below the day-ahead price?
\end{exercise}

\begin{exercise}[$\star$]\label{exo:s2:power-and-gas-trading:2}
A battery buys 2~MWh of storage at \euro0 and sells it at \euro100 with a round-trip efficiency of 88\%. What does it earn?
\end{exercise}

\begin{exercise}[$\star$]\label{exo:s2:power-and-gas-trading:3}
The forecast trade earns \euro2.54 per MWh on 61\% of hours with 1~MW. Check the annual figure over two years of 727 days.
\end{exercise}

\begin{exercise}[$\star\star$]\label{exo:s2:power-and-gas-trading:4}
Why does the forecast trade's edge fall so fast with a half-spread of \euro3?
\end{exercise}

\begin{exercise}[$\star\star$]\label{exo:s2:power-and-gas-trading:5}
Why are German prices most often negative at midday?
\end{exercise}

\begin{exercise}[$\star\star$]\label{exo:s2:power-and-gas-trading:6}
Kiesel and Paraschiv found an asymmetric response to forecast errors; Ziel found none significant. How could both be right?
\end{exercise}

\begin{exercise}[$\star\star\star$]\label{exo:s2:power-and-gas-trading:7}
\emph{Coding.} Run \texttt{battery\_table(PowerConfig(cycles=2))}. What does the second cycle add, and what does the table leave out?
\end{exercise}

\begin{exercise}[$\star\star\star$]\label{exo:s2:power-and-gas-trading:8}
\emph{Find the flaw.} ``Our battery backtest earned \euro69\,900 per MW a year on German day-ahead prices; that is what we will earn.''
\end{exercise}

\section{Problem: The Wind Forecast}

\begin{problem}[{Weekend problem --- power and gas trading}]\label{pb:s2:power-and-gas-trading:1}
Germany's day-ahead prices, the synthetic intraday market and the public record.

\medskip\noindent\textbf{Part I --- The markets.}
\begin{enumerate}
\item Define the \omterm{def:s2:power-and-gas-trading:daid}{day-ahead-intraday spread}.
\item Give the slope and correlation of price on residual load.
\item How are the synthetic intraday prices built, and what is planted?
\item What did Kiesel and Paraschiv and Ziel find?
\end{enumerate}

\medskip\noindent\textbf{Part II --- The forecast trade.}
\begin{enumerate}[resume]
\item Describe the trade.
\item Give its results at each skill and cost.
\item Define a \omterm{def:s2:power-and-gas-trading:weather}{weather-driven position}.
\item How does the trade die?
\end{enumerate}

\medskip\noindent\textbf{Part III --- The battery.}
\begin{enumerate}[resume]
\item Define \omterm{def:s2:power-and-gas-trading:battery}{battery arbitrage}.
\item Describe the price profile and negative hours.
\item How is the battery scheduled?
\item Give the battery table.
\end{enumerate}

\medskip\noindent\textbf{Part IV --- The verdict.}
\begin{enumerate}[resume]
\item State the \emph{named result}: the battery's revenue in day-ahead only and with intraday re-optimisation.
\item Why are both columns generous?
\item Why is the intraday extra small here?
\item What are the risks of a power book?
\item How would you backtest the battery honestly?
\item Which strategy file dies first as batteries multiply?
\item How does the forecast trade relate to chapter~15's flow trades?
\item In one sentence: what does a power trader sell?
\end{enumerate}
\end{problem}

\section{Interview questions}

\begin{interviewq}[$\star$ \iqroles{trader}]\label{iq:s2:power-and-gas-trading:1}
Why can power prices be negative?
\end{interviewq}

\begin{interviewq}[$\star\star$ \iqroles{researcher}]\label{iq:s2:power-and-gas-trading:2}
How would you forecast the \omterm{def:s2:power-and-gas-trading:daid}{day-ahead-intraday spread} for tomorrow's hours?
\end{interviewq}

\begin{interviewq}[$\star\star$ \iqroles{trader}]\label{iq:s2:power-and-gas-trading:3}
Your wind forecast is 3~GW above the market's for tomorrow afternoon. What do you do, and how much?
\end{interviewq}

\begin{interviewq}[$\star\star$ \iqroles{risk}]\label{iq:s2:power-and-gas-trading:4}
How would you set limits for an intraday power desk?
\end{interviewq}

\begin{interviewq}[$\star\star$ \iqroles{developer}]\label{iq:s2:power-and-gas-trading:5}
Design the system that schedules a fleet of batteries every quarter hour.
\end{interviewq}

\begin{interviewq}[$\star\star\star$ \iqroles{researcher}]\label{iq:s2:power-and-gas-trading:6}
A battery with round-trip efficiency $\eta$ buys at $p_1$ and sells at $p_2$. Show that the trade pays only if $p_2 > p_1/\eta$ (for positive
prices), and say what changes when $p_1 < 0$.
\end{interviewq}
